% Options for packages loaded elsewhere
\PassOptionsToPackage{unicode}{hyperref}
\PassOptionsToPackage{hyphens}{url}
%
\documentclass[
  24pt,
  ignorenonframetext,
  aspectratio=43,
]{beamer}
\usepackage{pgfpages}
\setbeamertemplate{caption}[numbered]
\setbeamertemplate{caption label separator}{: }
\setbeamercolor{caption name}{fg=normal text.fg}
\beamertemplatenavigationsymbolsvertical
% Prevent slide breaks in the middle of a paragraph
\widowpenalties 1 10000
\raggedbottom
\setbeamertemplate{part page}{
  \centering
  \begin{beamercolorbox}[sep=16pt,center]{part title}
    \usebeamerfont{part title}\insertpart\par
  \end{beamercolorbox}
}
\setbeamertemplate{section page}{
  \centering
  \begin{beamercolorbox}[sep=12pt,center]{section title}
    \usebeamerfont{section title}\insertsection\par
  \end{beamercolorbox}
}
\setbeamertemplate{subsection page}{
  \centering
  \begin{beamercolorbox}[sep=8pt,center]{subsection title}
    \usebeamerfont{subsection title}\insertsubsection\par
  \end{beamercolorbox}
}
\AtBeginPart{
  \frame{\partpage}
}
\AtBeginSection{
  \ifbibliography
  \else
    \frame{\sectionpage}
  \fi
}
\AtBeginSubsection{
  \frame{\subsectionpage}
}

\usepackage{amsmath,amssymb}
\usepackage{iftex}
\ifPDFTeX
  \usepackage[T1]{fontenc}
  \usepackage[utf8]{inputenc}
  \usepackage{textcomp} % provide euro and other symbols
\else % if luatex or xetex
  \usepackage{unicode-math}
  \defaultfontfeatures{Scale=MatchLowercase}
  \defaultfontfeatures[\rmfamily]{Ligatures=TeX,Scale=1}
\fi
\usepackage{lmodern}
\usetheme[]{CambridgeUS}
\usecolortheme{beaver}
\ifPDFTeX\else  
    % xetex/luatex font selection
\fi
% Use upquote if available, for straight quotes in verbatim environments
\IfFileExists{upquote.sty}{\usepackage{upquote}}{}
\IfFileExists{microtype.sty}{% use microtype if available
  \usepackage[]{microtype}
  \UseMicrotypeSet[protrusion]{basicmath} % disable protrusion for tt fonts
}{}
\makeatletter
\@ifundefined{KOMAClassName}{% if non-KOMA class
  \IfFileExists{parskip.sty}{%
    \usepackage{parskip}
  }{% else
    \setlength{\parindent}{0pt}
    \setlength{\parskip}{6pt plus 2pt minus 1pt}}
}{% if KOMA class
  \KOMAoptions{parskip=half}}
\makeatother
\usepackage{xcolor}
\newif\ifbibliography
\setlength{\emergencystretch}{3em} % prevent overfull lines
\setcounter{secnumdepth}{-\maxdimen} % remove section numbering

\usepackage{color}
\usepackage{fancyvrb}
\newcommand{\VerbBar}{|}
\newcommand{\VERB}{\Verb[commandchars=\\\{\}]}
\DefineVerbatimEnvironment{Highlighting}{Verbatim}{commandchars=\\\{\}}
% Add ',fontsize=\small' for more characters per line
\usepackage{framed}
\definecolor{shadecolor}{RGB}{241,243,245}
\newenvironment{Shaded}{\begin{snugshade}}{\end{snugshade}}
\newcommand{\AlertTok}[1]{\textcolor[rgb]{0.68,0.00,0.00}{#1}}
\newcommand{\AnnotationTok}[1]{\textcolor[rgb]{0.37,0.37,0.37}{#1}}
\newcommand{\AttributeTok}[1]{\textcolor[rgb]{0.40,0.45,0.13}{#1}}
\newcommand{\BaseNTok}[1]{\textcolor[rgb]{0.68,0.00,0.00}{#1}}
\newcommand{\BuiltInTok}[1]{\textcolor[rgb]{0.00,0.23,0.31}{#1}}
\newcommand{\CharTok}[1]{\textcolor[rgb]{0.13,0.47,0.30}{#1}}
\newcommand{\CommentTok}[1]{\textcolor[rgb]{0.37,0.37,0.37}{#1}}
\newcommand{\CommentVarTok}[1]{\textcolor[rgb]{0.37,0.37,0.37}{\textit{#1}}}
\newcommand{\ConstantTok}[1]{\textcolor[rgb]{0.56,0.35,0.01}{#1}}
\newcommand{\ControlFlowTok}[1]{\textcolor[rgb]{0.00,0.23,0.31}{\textbf{#1}}}
\newcommand{\DataTypeTok}[1]{\textcolor[rgb]{0.68,0.00,0.00}{#1}}
\newcommand{\DecValTok}[1]{\textcolor[rgb]{0.68,0.00,0.00}{#1}}
\newcommand{\DocumentationTok}[1]{\textcolor[rgb]{0.37,0.37,0.37}{\textit{#1}}}
\newcommand{\ErrorTok}[1]{\textcolor[rgb]{0.68,0.00,0.00}{#1}}
\newcommand{\ExtensionTok}[1]{\textcolor[rgb]{0.00,0.23,0.31}{#1}}
\newcommand{\FloatTok}[1]{\textcolor[rgb]{0.68,0.00,0.00}{#1}}
\newcommand{\FunctionTok}[1]{\textcolor[rgb]{0.28,0.35,0.67}{#1}}
\newcommand{\ImportTok}[1]{\textcolor[rgb]{0.00,0.46,0.62}{#1}}
\newcommand{\InformationTok}[1]{\textcolor[rgb]{0.37,0.37,0.37}{#1}}
\newcommand{\KeywordTok}[1]{\textcolor[rgb]{0.00,0.23,0.31}{\textbf{#1}}}
\newcommand{\NormalTok}[1]{\textcolor[rgb]{0.00,0.23,0.31}{#1}}
\newcommand{\OperatorTok}[1]{\textcolor[rgb]{0.37,0.37,0.37}{#1}}
\newcommand{\OtherTok}[1]{\textcolor[rgb]{0.00,0.23,0.31}{#1}}
\newcommand{\PreprocessorTok}[1]{\textcolor[rgb]{0.68,0.00,0.00}{#1}}
\newcommand{\RegionMarkerTok}[1]{\textcolor[rgb]{0.00,0.23,0.31}{#1}}
\newcommand{\SpecialCharTok}[1]{\textcolor[rgb]{0.37,0.37,0.37}{#1}}
\newcommand{\SpecialStringTok}[1]{\textcolor[rgb]{0.13,0.47,0.30}{#1}}
\newcommand{\StringTok}[1]{\textcolor[rgb]{0.13,0.47,0.30}{#1}}
\newcommand{\VariableTok}[1]{\textcolor[rgb]{0.07,0.07,0.07}{#1}}
\newcommand{\VerbatimStringTok}[1]{\textcolor[rgb]{0.13,0.47,0.30}{#1}}
\newcommand{\WarningTok}[1]{\textcolor[rgb]{0.37,0.37,0.37}{\textit{#1}}}

\providecommand{\tightlist}{%
  \setlength{\itemsep}{0pt}\setlength{\parskip}{0pt}}\usepackage{longtable,booktabs,array}
\usepackage{calc} % for calculating minipage widths
\usepackage{caption}
% Make caption package work with longtable
\makeatletter
\def\fnum@table{\tablename~\thetable}
\makeatother
\usepackage{graphicx}
\makeatletter
\newsavebox\pandoc@box
\newcommand*\pandocbounded[1]{% scales image to fit in text height/width
  \sbox\pandoc@box{#1}%
  \Gscale@div\@tempa{\textheight}{\dimexpr\ht\pandoc@box+\dp\pandoc@box\relax}%
  \Gscale@div\@tempb{\linewidth}{\wd\pandoc@box}%
  \ifdim\@tempb\p@<\@tempa\p@\let\@tempa\@tempb\fi% select the smaller of both
  \ifdim\@tempa\p@<\p@\scalebox{\@tempa}{\usebox\pandoc@box}%
  \else\usebox{\pandoc@box}%
  \fi%
}
% Set default figure placement to htbp
\def\fps@figure{htbp}
\makeatother
% definitions for citeproc citations
\NewDocumentCommand\citeproctext{}{}
\NewDocumentCommand\citeproc{mm}{%
  \begingroup\def\citeproctext{#2}\cite{#1}\endgroup}
\makeatletter
 % allow citations to break across lines
 \let\@cite@ofmt\@firstofone
 % avoid brackets around text for \cite:
 \def\@biblabel#1{}
 \def\@cite#1#2{{#1\if@tempswa , #2\fi}}
\makeatother
\newlength{\cslhangindent}
\setlength{\cslhangindent}{1.5em}
\newlength{\csllabelwidth}
\setlength{\csllabelwidth}{3em}
\newenvironment{CSLReferences}[2] % #1 hanging-indent, #2 entry-spacing
 {\begin{list}{}{%
  \setlength{\itemindent}{0pt}
  \setlength{\leftmargin}{0pt}
  \setlength{\parsep}{0pt}
  % turn on hanging indent if param 1 is 1
  \ifodd #1
   \setlength{\leftmargin}{\cslhangindent}
   \setlength{\itemindent}{-1\cslhangindent}
  \fi
  % set entry spacing
  \setlength{\itemsep}{#2\baselineskip}}}
 {\end{list}}
\usepackage{calc}
\newcommand{\CSLBlock}[1]{\hfill\break\parbox[t]{\linewidth}{\strut\ignorespaces#1\strut}}
\newcommand{\CSLLeftMargin}[1]{\parbox[t]{\csllabelwidth}{\strut#1\strut}}
\newcommand{\CSLRightInline}[1]{\parbox[t]{\linewidth - \csllabelwidth}{\strut#1\strut}}
\newcommand{\CSLIndent}[1]{\hspace{\cslhangindent}#1}

\makeatletter
\@ifpackageloaded{caption}{}{\usepackage{caption}}
\AtBeginDocument{%
\ifdefined\contentsname
  \renewcommand*\contentsname{Table of contents}
\else
  \newcommand\contentsname{Table of contents}
\fi
\ifdefined\listfigurename
  \renewcommand*\listfigurename{List of Figures}
\else
  \newcommand\listfigurename{List of Figures}
\fi
\ifdefined\listtablename
  \renewcommand*\listtablename{List of Tables}
\else
  \newcommand\listtablename{List of Tables}
\fi
\ifdefined\figurename
  \renewcommand*\figurename{Figure}
\else
  \newcommand\figurename{Figure}
\fi
\ifdefined\tablename
  \renewcommand*\tablename{Table}
\else
  \newcommand\tablename{Table}
\fi
}
\@ifpackageloaded{float}{}{\usepackage{float}}
\floatstyle{ruled}
\@ifundefined{c@chapter}{\newfloat{codelisting}{h}{lop}}{\newfloat{codelisting}{h}{lop}[chapter]}
\floatname{codelisting}{Listing}
\newcommand*\listoflistings{\listof{codelisting}{List of Listings}}
\makeatother
\makeatletter
\makeatother
\makeatletter
\@ifpackageloaded{caption}{}{\usepackage{caption}}
\@ifpackageloaded{subcaption}{}{\usepackage{subcaption}}
\makeatother

\usepackage{bookmark}

\IfFileExists{xurl.sty}{\usepackage{xurl}}{} % add URL line breaks if available
\urlstyle{same} % disable monospaced font for URLs
\hypersetup{
  pdftitle={Fast Robust Moments with fromo},
  pdfauthor={Steven E. Pav},
  hidelinks,
  pdfcreator={LaTeX via pandoc}}


\title{Fast Robust Moments with \texttt{fromo}}
\author{Steven E. Pav}
\date{}

\begin{document}
\frame{\titlepage}


\begin{frame}[fragile]{\texttt{fromo} package}
\phantomsection\label{fromo-package}
\footnotesize

\textbf{R} package to compute Fast Robust Moments.\\
Get from CRAN:
\texttt{install.packages(c(\textquotesingle{}fromo\textquotesingle{}))}.

\textbf{Fast}

\begin{itemize}
\tightlist
\item
  Coded in \texttt{Rcpp}, heavy use of templating optimizations.
\item
  Running computations add and subtract observations, rather than
  compute entirely from scratch.
\end{itemize}

\textbf{(Numerically) Robust}

\begin{itemize}
\item
  Uses Welford-Terriberry method to avoid disastrous cancellation.
  (Welford 1962; Terriberry 2008; Pébay et al. 2016)
\item
  Checks for impossible outcomes (negative even moments) due to
  numerical issues and recomputes.
\end{itemize}

\textbf{Moments}

\begin{itemize}
\tightlist
\item
  Computes univariate sums, means, centered moments, cumulants;
  bivariate covariance, correlation and regression.
\end{itemize}
\end{frame}

\begin{frame}{Running Computations I}
\phantomsection\label{running-computations-i}
Running (``rolling'') computations performed over a window of size
\(w\).

\begin{center}
\includegraphics[width=0.7\linewidth,height=\textheight,keepaspectratio]{g55079.png}
\end{center}

To move the window forward, add \(x_n\) and remove \(x_{n-w}\) from the
sum.

This is much more efficient than computing from scratch.
\end{frame}

\begin{frame}{Time Running Computations I}
\phantomsection\label{time-running-computations-i}
The time running computations are over a time window of size \(w\).

\begin{itemize}
\tightlist
\item
  The data \(x_i\) have associated times \(t_i\).
\item
  Moments are computed as of lookback times \(lb_k\):

  \begin{itemize}
  \tightlist
  \item
    Includes \(x_i\) such that \(lb_k - w < t_i \le lb_k\).
  \end{itemize}
\item
  Lookback times default to the \(t_i\).
\end{itemize}

\begin{center}
\includegraphics[width=0.7\linewidth,height=\textheight,keepaspectratio]{g152441a.png}
\end{center}
\end{frame}

\begin{frame}{Time Running Computations II}
\phantomsection\label{time-running-computations-ii}
\begin{center}
\includegraphics[width=0.7\linewidth,height=\textheight,keepaspectratio]{g152441a.png}
\end{center}

To move the window forward, add \(x_j\) and remove \(x_i, x_{i+1}\) from
the sum.

\begin{center}
\includegraphics[width=0.7\linewidth,height=\textheight,keepaspectratio]{g152441b.png}
\end{center}
\end{frame}

\begin{frame}[fragile]{Other Features}
\phantomsection\label{other-features}
\begin{itemize}
\tightlist
\item
  Static (non-running) moment objects support algebraic operations.
\item
  All computations support optional weights.

  \begin{itemize}
  \tightlist
  \item
    Replication weights by default, but code can auto-normalize to mean
    1.
  \end{itemize}
\item
  Computations will ignore \texttt{NA} or propagate them based on a
  flag.

  \begin{itemize}
  \tightlist
  \item
    Code is faster without \texttt{NA} checking.
  \end{itemize}
\item
  Running computations have tweakable ``recompute counter'' to trigger
  compute from scratch.

  \begin{itemize}
  \tightlist
  \item
    Fine tune between speed and robustness.
  \end{itemize}
\item
  Time running computations support variable width window.

  \begin{itemize}
  \tightlist
  \item
    Computes over all observations since last lookback time.
  \end{itemize}
\item
  Time running computations can use cumulative weight as time.
\item
  Running computations support adjustments like Z-scoring or rescaling.

  \begin{itemize}
  \tightlist
  \item
    These support lookahead/lookback to adjust to moments over a leading
    or lagging window.
  \end{itemize}
\item
  Bivariate (time) running computations supported for aligned \(x, y\).
\end{itemize}
\end{frame}

\begin{frame}[fragile]{Usage I}
\phantomsection\label{usage-i}
Input data are daily Market Returns. (French 2020)

\AddToHookNext{env/Highlighting/begin}{\tiny}

\begin{Shaded}
\begin{Highlighting}[numbers=left,,]
\CommentTok{\# remotes::install\_github("shabbychef/tsrsa/rpkg")}
\FunctionTok{library}\NormalTok{(tsrsa)}
\CommentTok{\# daily \textquotesingle{}returns\textquotesingle{} of Fama French 4 factors}
\FunctionTok{data}\NormalTok{(dff4)}
\NormalTok{mkt }\OtherTok{\textless{}{-}}\NormalTok{ dff4}\SpecialCharTok{$}\NormalTok{Mkt}
\NormalTok{smb }\OtherTok{\textless{}{-}}\NormalTok{ dff4}\SpecialCharTok{$}\NormalTok{SMB}
\FunctionTok{plot}\NormalTok{(mkt)}
\end{Highlighting}
\end{Shaded}

\begin{figure}[H]

{\centering \includegraphics[width=1\linewidth,height=0.4\textheight]{fromo_files/figure-beamer/loadem-1.pdf}

}

\caption{The daily percent returns of the Fama French Market Factor.}

\end{figure}%
\end{frame}

\begin{frame}[fragile]{Usage II}
\phantomsection\label{usage-ii}
Compute the running yearly mean returns.

\AddToHookNext{env/Highlighting/begin}{\tiny}

\begin{Shaded}
\begin{Highlighting}[numbers=left,,]
\FunctionTok{library}\NormalTok{(fromo)}
\NormalTok{mu\_mkt }\OtherTok{\textless{}{-}} \FunctionTok{running\_mean}\NormalTok{(mkt,}\AttributeTok{win=}\DecValTok{253}\NormalTok{,}\AttributeTok{min\_df=}\DecValTok{30}\NormalTok{)}
\NormalTok{mu\_mkt }\OtherTok{\textless{}{-}} \FunctionTok{setNames}\NormalTok{(}\FunctionTok{.xts}\NormalTok{(mu\_mkt,}\AttributeTok{index=}\FunctionTok{index}\NormalTok{(mkt)),}\StringTok{\textquotesingle{}mu\_Mkt\textquotesingle{}}\NormalTok{)}
\FunctionTok{plot}\NormalTok{(mu\_mkt)}
\end{Highlighting}
\end{Shaded}

\begin{figure}[H]

{\centering \includegraphics[width=1\linewidth,height=0.4\textheight]{fromo_files/figure-beamer/usage_II-1.pdf}

}

\caption{The running yearly average percent returns of the Market
Factor.}

\end{figure}%
\end{frame}

\begin{frame}[fragile]{Usage III}
\phantomsection\label{usage-iii}
Compute the three month (63 day) running stdev of returns.

\AddToHookNext{env/Highlighting/begin}{\tiny}

\begin{Shaded}
\begin{Highlighting}[numbers=left,,]
\NormalTok{mkt\_sd }\OtherTok{\textless{}{-}} \FunctionTok{running\_sd3}\NormalTok{(mkt,}\AttributeTok{win=}\DecValTok{63}\NormalTok{,}\AttributeTok{min\_df=}\DecValTok{30}\NormalTok{)}
\NormalTok{mkt\_sd }\OtherTok{\textless{}{-}} \FunctionTok{setNames}\NormalTok{(}\FunctionTok{.xts}\NormalTok{(mkt\_sd[,}\DecValTok{1}\NormalTok{],}\AttributeTok{index=}\FunctionTok{index}\NormalTok{(mkt)),}\StringTok{\textquotesingle{}skew\_Mkt\textquotesingle{}}\NormalTok{)}
\FunctionTok{plot}\NormalTok{(mkt\_sd)}
\end{Highlighting}
\end{Shaded}

\begin{figure}[H]

{\centering \includegraphics[width=1\linewidth,height=0.4\textheight]{fromo_files/figure-beamer/usage_III-1.pdf}

}

\caption{The running 63 day standard devation of percent returns of the
Market Factor.}

\end{figure}%
\end{frame}

\begin{frame}[fragile]{Usage IV}
\phantomsection\label{usage-iv}
Compute the running yearly skew of returns.
\AddToHookNext{env/Highlighting/begin}{\tiny}

\begin{Shaded}
\begin{Highlighting}[numbers=left,,]
\NormalTok{mu\_sku }\OtherTok{\textless{}{-}} \FunctionTok{running\_skew4}\NormalTok{(mkt,}\AttributeTok{win=}\DecValTok{253}\NormalTok{,}\AttributeTok{min\_df=}\DecValTok{30}\NormalTok{)}
\NormalTok{mu\_sku }\OtherTok{\textless{}{-}} \FunctionTok{setNames}\NormalTok{(}\FunctionTok{.xts}\NormalTok{(mu\_sku[,}\DecValTok{1}\NormalTok{],}\AttributeTok{index=}\FunctionTok{index}\NormalTok{(mkt)),}\StringTok{\textquotesingle{}skew\_Mkt\textquotesingle{}}\NormalTok{)}
\FunctionTok{plot}\NormalTok{(mu\_sku)}
\end{Highlighting}
\end{Shaded}

\begin{figure}[H]

{\centering \includegraphics[width=1\linewidth,height=0.4\textheight]{fromo_files/figure-beamer/usage_IV-1.pdf}

}

\caption{The running yearly average percent returns of the Market
Factor.}

\end{figure}%
\end{frame}

\begin{frame}[fragile]{Usage V}
\phantomsection\label{usage-v}
Z-score the returns, with a lookback window of 10 trading days.

\AddToHookNext{env/Highlighting/begin}{\tiny}

\begin{Shaded}
\begin{Highlighting}[numbers=left,,]
\NormalTok{z\_score }\OtherTok{\textless{}{-}} \FunctionTok{running\_zscored}\NormalTok{(mkt,}\AttributeTok{win=}\DecValTok{253}\NormalTok{,}\AttributeTok{lookahead=}\SpecialCharTok{{-}}\DecValTok{10}\NormalTok{,}\AttributeTok{min\_df=}\DecValTok{30}\NormalTok{)}
\NormalTok{z\_score }\OtherTok{\textless{}{-}} \FunctionTok{setNames}\NormalTok{(}\FunctionTok{.xts}\NormalTok{(z\_score,}\AttributeTok{index=}\FunctionTok{index}\NormalTok{(mkt)),}\StringTok{\textquotesingle{}Z of Mkt\textquotesingle{}}\NormalTok{)}
\FunctionTok{plot}\NormalTok{(z\_score)}
\end{Highlighting}
\end{Shaded}

\begin{figure}[H]

{\centering \includegraphics[width=1\linewidth,height=0.4\textheight]{fromo_files/figure-beamer/usage_V-1.pdf}

}

\caption{Z-scores of the daily percent returns of the Market Factor.}

\end{figure}%
\end{frame}

\begin{frame}[fragile]{Usage VI}
\phantomsection\label{usage-vi}
Compute trailing 3 month total return, using time-running computation.

\AddToHookNext{env/Highlighting/begin}{\tiny}

\begin{Shaded}
\begin{Highlighting}[numbers=left,,]
\CommentTok{\# compute on log returns, which telescope, then convert back to percents.}
\NormalTok{t\_sum }\OtherTok{\textless{}{-}} \FunctionTok{t\_running\_sum}\NormalTok{(}\FunctionTok{log}\NormalTok{(}\DecValTok{1}\FloatTok{+0.01}\SpecialCharTok{*}\NormalTok{mkt),}\AttributeTok{time=}\FunctionTok{index}\NormalTok{(mkt),}\AttributeTok{window=}\DecValTok{365}\SpecialCharTok{/}\DecValTok{4}\NormalTok{,}\AttributeTok{min\_df=}\DecValTok{40}\NormalTok{)}
\NormalTok{t\_sum }\OtherTok{\textless{}{-}} \DecValTok{100}\SpecialCharTok{*}\NormalTok{(}\FunctionTok{exp}\NormalTok{(t\_sum) }\SpecialCharTok{{-}} \DecValTok{1}\NormalTok{)}
\NormalTok{t\_sum }\OtherTok{\textless{}{-}} \FunctionTok{setNames}\NormalTok{(}\FunctionTok{.xts}\NormalTok{(t\_sum,}\AttributeTok{index=}\FunctionTok{index}\NormalTok{(mkt)),}\StringTok{\textquotesingle{}trailing total returns of Mkt\textquotesingle{}}\NormalTok{)}
\FunctionTok{plot}\NormalTok{(t\_sum)}
\end{Highlighting}
\end{Shaded}

\begin{figure}[H]

{\centering \includegraphics[width=1\linewidth,height=0.4\textheight]{fromo_files/figure-beamer/usage_VI-1.pdf}

}

\caption{Total returns over previous 3 months of the Market Factor.}

\end{figure}%
\end{frame}

\begin{frame}[fragile]{Usage VII}
\phantomsection\label{usage-vii}
Compute trailing quarterly standard deviation of returns, with
non-overlapping windows.

\AddToHookNext{env/Highlighting/begin}{\tiny}

\begin{Shaded}
\begin{Highlighting}[numbers=left,,]
\NormalTok{times }\OtherTok{\textless{}{-}} \FunctionTok{index}\NormalTok{(mkt)}
\FunctionTok{library}\NormalTok{(lubridate)}
\NormalTok{qt\_ends }\OtherTok{\textless{}{-}} \FunctionTok{unique}\NormalTok{(lubridate}\SpecialCharTok{::}\FunctionTok{ceiling\_date}\NormalTok{(times,}\StringTok{\textquotesingle{}quarter\textquotesingle{}}\NormalTok{) }\SpecialCharTok{\%m{-}\%} \FunctionTok{days}\NormalTok{(}\DecValTok{1}\NormalTok{))}
\NormalTok{t\_sd }\OtherTok{\textless{}{-}} \FunctionTok{t\_running\_sd}\NormalTok{(mkt,}\AttributeTok{time=}\NormalTok{times,}\AttributeTok{lb\_time=}\NormalTok{qt\_ends,}\AttributeTok{window=}\ConstantTok{NULL}\NormalTok{,}\AttributeTok{variable\_win=}\ConstantTok{TRUE}\NormalTok{) }
\NormalTok{t\_sd }\OtherTok{\textless{}{-}} \FunctionTok{setNames}\NormalTok{(}\FunctionTok{.xts}\NormalTok{(t\_sd,}\AttributeTok{index=}\NormalTok{qt\_ends),}\StringTok{\textquotesingle{}trailing quarterly vol, returns of Mkt\textquotesingle{}}\NormalTok{)}
\FunctionTok{plot}\NormalTok{(t\_sd)}
\end{Highlighting}
\end{Shaded}

\begin{figure}[H]

{\centering \includegraphics[width=1\linewidth,height=0.4\textheight]{fromo_files/figure-beamer/usage_VII-1.pdf}

}

\caption{Non-overlapping quarterly stdev of the Market Factor.}

\end{figure}%
\end{frame}

\begin{frame}[fragile]{Usage VIII}
\phantomsection\label{usage-viii}
Compute beta of SMB vs Mkt on quarterly basis.

\AddToHookNext{env/Highlighting/begin}{\tiny}

\begin{Shaded}
\begin{Highlighting}[numbers=left,,]
\NormalTok{times }\OtherTok{\textless{}{-}} \FunctionTok{index}\NormalTok{(mkt)}
\FunctionTok{library}\NormalTok{(lubridate)}
\NormalTok{qt\_ends }\OtherTok{\textless{}{-}} \FunctionTok{unique}\NormalTok{(lubridate}\SpecialCharTok{::}\FunctionTok{ceiling\_date}\NormalTok{(times,}\StringTok{\textquotesingle{}quarter\textquotesingle{}}\NormalTok{) }\SpecialCharTok{\%m{-}\%} \FunctionTok{days}\NormalTok{(}\DecValTok{1}\NormalTok{))}
\NormalTok{beta }\OtherTok{\textless{}{-}} \FunctionTok{t\_running\_regression\_slope}\NormalTok{(}\AttributeTok{x=}\NormalTok{mkt,}\AttributeTok{y=}\NormalTok{smb,}\AttributeTok{time=}\NormalTok{times,}\AttributeTok{lb\_time=}\NormalTok{qt\_ends,}\AttributeTok{window=}\ConstantTok{NULL}\NormalTok{,}\AttributeTok{variable\_win=}\ConstantTok{TRUE}\NormalTok{) }
\NormalTok{beta }\OtherTok{\textless{}{-}} \FunctionTok{setNames}\NormalTok{(}\FunctionTok{.xts}\NormalTok{(beta,}\AttributeTok{index=}\NormalTok{qt\_ends),}\StringTok{\textquotesingle{}trailing quarterly beta, SMB vs Mkt\textquotesingle{}}\NormalTok{)}
\FunctionTok{plot}\NormalTok{(beta)}
\end{Highlighting}
\end{Shaded}

\begin{figure}[H]

{\centering \includegraphics[width=1\linewidth,height=0.4\textheight]{fromo_files/figure-beamer/usage_VIII-1.pdf}

}

\caption{Non-overlapping quarterly beta of SMB against the Market
Factor.}

\end{figure}%
\end{frame}

\begin{frame}[fragile]{Usage IX}
\phantomsection\label{usage-ix}
Compute correlation of SMB vs Mkt on quarterly basis, upweight Fridays
by 4x. (Just Because.)

\AddToHookNext{env/Highlighting/begin}{\tiny}

\begin{Shaded}
\begin{Highlighting}[numbers=left,,]
\NormalTok{times }\OtherTok{\textless{}{-}} \FunctionTok{index}\NormalTok{(mkt)}
\FunctionTok{library}\NormalTok{(lubridate)}
\NormalTok{wts }\OtherTok{\textless{}{-}} \DecValTok{1} \SpecialCharTok{+} \DecValTok{3} \SpecialCharTok{*}\NormalTok{ (}\FunctionTok{weekdays}\NormalTok{(times)}\SpecialCharTok{==}\StringTok{\textquotesingle{}Friday\textquotesingle{}}\NormalTok{)}
\NormalTok{qt\_ends }\OtherTok{\textless{}{-}} \FunctionTok{unique}\NormalTok{(lubridate}\SpecialCharTok{::}\FunctionTok{ceiling\_date}\NormalTok{(times,}\StringTok{\textquotesingle{}quarter\textquotesingle{}}\NormalTok{) }\SpecialCharTok{\%m{-}\%} \FunctionTok{days}\NormalTok{(}\DecValTok{1}\NormalTok{))}
\NormalTok{rho }\OtherTok{\textless{}{-}} \FunctionTok{t\_running\_correlation}\NormalTok{(}\AttributeTok{x=}\NormalTok{mkt,}\AttributeTok{y=}\NormalTok{smb,}\AttributeTok{time=}\NormalTok{times,}\AttributeTok{wts=}\NormalTok{wts,}\AttributeTok{lb\_time=}\NormalTok{qt\_ends,}\AttributeTok{window=}\ConstantTok{NULL}\NormalTok{,}\AttributeTok{variable\_win=}\ConstantTok{TRUE}\NormalTok{) }
\NormalTok{rho }\OtherTok{\textless{}{-}} \FunctionTok{setNames}\NormalTok{(}\FunctionTok{.xts}\NormalTok{(rho,}\AttributeTok{index=}\NormalTok{qt\_ends),}\StringTok{\textquotesingle{}trailing weighted quarterly rho, SMB vs Mkt\textquotesingle{}}\NormalTok{)}
\FunctionTok{plot}\NormalTok{(rho)}
\end{Highlighting}
\end{Shaded}

\begin{figure}[H]

{\centering \includegraphics[width=1\linewidth,height=0.4\textheight]{fromo_files/figure-beamer/usage_IX-1.pdf}

}

\caption{Non-overlapping quarterly correlation of SMB against the
Market.}

\end{figure}%
\end{frame}

\begin{frame}{References}
\phantomsection\label{references}
\footnotesize

\phantomsection\label{refs}
\begin{CSLReferences}{1}{0}
\bibitem[\citeproctext]{ref-french_data_library}
French, Kenneth. 2020. {``Data Library.''} Privately Published.
\url{http://mba.tuck.dartmouth.edu/pages/faculty/ken.french/data_library.html}.

\bibitem[\citeproctext]{ref-fromopkg}
Pav, Steven E. 2024. \emph{{fromo}: Fast Robust Moments}.
\url{https://github.com/shabbychef/fromo}.

\bibitem[\citeproctext]{ref-pebay2016numerically}
Pébay, Philippe, Timothy B Terriberry, Hemanth Kolla, and Janine
Bennett. 2016. {``Numerically Stable, Scalable Formulas for Parallel and
Online Computation of Higher-Order Multivariate Central Moments with
Arbitrary Weights.''} \emph{Computational Statistics} 31: 1305--25.
\url{https://www.osti.gov/biblio/1426900}.

\bibitem[\citeproctext]{ref-terriberry2008computing}
Terriberry, Timothy B. 2008. {``Computing Higher-Order Moments
Online.''}
\url{https://web.archive.org/web/20140423031833/http://people.xiph.org/~tterribe/notes/homs.html}.

\bibitem[\citeproctext]{ref-welford1962note}
Welford, Barry Payne. 1962. {``Note on a Method for Calculating
Corrected Sums of Squares and Products.''} \emph{Technometrics} 4 (3):
419--20. \url{https://doi.org/10.1080/00401706.1962.10490022}.

\end{CSLReferences}
\end{frame}




\end{document}
